% ------------------------------------------------------------------------------
% Metodologia
% ------------------------------------------------------------------------------
\chapter{Metodologia}
\label{chap_metodologia}
Este capítulo descreve os procedimentos adotados para a realização dos experimentos computacionais de queda livre, detalhando os materiais e ferramentas utilizados, os parâmetros de entrada considerados, o método de simulação empregado e a abordagem para análise dos dados de acurácia.

\section{Materiais Computacionais e Organização dos Dados}


As simulações foram realizadas em ambiente Python, utilizando as bibliotecas \texttt{numpy} para cálculos numéricos, \texttt{pandas} para manipulação de dados e \texttt{matplotlib} para geração de gráficos. Todo o código foi estruturado em um notebook Jupyter, permitindo a execução sequencial dos experimentos e a documentação dos resultados.

O código-fonte completo está disponível publicamente em \cite{github_queda_livre}.

Os experimentos computacionais foram realizados em três etapas sequenciais, cada uma abordando um aspecto físico distinto do movimento:
    \begin{enumerate}
        \item \textbf{Queda livre de partícula:} Simulação do movimento de uma partícula ideal em queda livre, desconsiderando qualquer tipo de resistência do ar. Para partículas pequenas, densas e em trajetórias de curta distância, a resistência do ar é desprezível em relação ao peso, conforme discutido em diversas fontes da literatura (por exemplo, Halliday, Resnick e Tipler)\cite{halliday2011,tipler2009}. Por isso, este modelo é amplamente utilizado como referência para validação de métodos numéricos e comparação com a solução analítica clássica.
        \item \textbf{Corpo rígido sem resistência do ar:} Análise do movimento considerando um corpo rígido, discutindo as diferenças conceituais e possíveis impactos na modelagem, ainda sem resistência do ar.
        \begin{enumerate}
        \item \textbf{Corpo rígido com resistência do ar linear:} Simulação do movimento de um corpo rígido considerando a presença de resistência do ar proporcional à velocidade (modelo linear, Lei de Stokes), permitindo avaliar o efeito desse fator sobre o movimento e a acurácia do método numérico.
        \item \textbf{Corpo rígido com resistência do ar quadrática:} Simulação do movimento de um corpo rígido considerando a resistência do ar proporcional ao quadrado da velocidade (modelo quadrático), possibilitando a análise do comportamento do sistema em regimes de velocidade mais elevados e a comparação entre os diferentes modelos de resistência.
        \end{enumerate}
        \item \textbf{Análise de acurácia dos métodos numéricos:} Para cada modelo físico, foi realizada uma análise quantitativa da acurácia dos métodos numéricos empregados, variando o passo de tempo ($\Delta t$) e comparando os resultados numéricos com as soluções analíticas de referência. Essa análise permitiu avaliar a convergência dos métodos e identificar as condições sob as quais a precisão é adequada para cada cenário.
    \end{enumerate}

\section{Procedimento Experimental}
\label{sec_procedimento_experimental}

Para cada caso, os parâmetros de entrada foram lidos dos arquivos de input e utilizados para configurar as simulações. O método de Euler foi empregado para resolver numericamente as equações diferenciais do movimento, permitindo o cálculo iterativo da posição, velocidade e tempo até o corpo atingir o solo.

A única variação sistemática realizada nos experimentos foi no passo de tempo ($\Delta t$), com o objetivo de avaliar a influência desse parâmetro na precisão dos resultados numéricos. Os demais parâmetros, incluindo os coeficientes de resistência do ar, permaneceram fixos conforme detalhado na seção de parâmetros principais.

Para cada simulação, foram analisados os resultados de posição, velocidade e tempo final, permitindo a comparação entre os diferentes modelos e a validação dos métodos numéricos empregados.

\subsection{Parâmetros de Entrada das Simulações}
\label{sec_parametros_entrada}

Os experimentos computacionais foram organizados de acordo com o modelo físico analisado, utilizando parâmetros realistas para garantir a representatividade dos cenários e permitir comparações rigorosas entre os regimes de resistência do ar.

Os valores adotados para cada cenário estão detalhados nas Tabelas~\ref{tab:parametros_queda_livre}, \ref{tab:parametros_linear} e \ref{tab:parametros_quadratica}, que apresentam, respectivamente, os parâmetros para a partícula sem resistência, para o corpo rígido com resistência linear e para o corpo rígido com resistência quadrática. Essas tabelas permitem rastreabilidade, comparação entre modelos e reprodutibilidade científica.

\textbf{Parâmetros utilizados nos experimentos:}
\begin{description}
    \item[Queda Livre de Partícula:] Altura inicial ($H$), velocidade inicial ($V$), aceleração da gravidade ($g$).\\
    \begin{table}[H]
\centering
\caption{Tabela – Dados referentes a Parametros queda livre.}
\label{tab:parametros_queda_livre}
\begin{tabular}{llll}
\toprule
Parâmetro & Descrição & Valor & Unidade \\
\midrule
H & Altura Inicial & 10000.0 & m \\
V & Velocidade Inicial & 0.0 & m/s \\
g & Aceleração da Gravidade & 9.81 & m/s² \\
\bottomrule
\end{tabular}

\end{table}

    \label{tab:parametros_queda_livre}
    \item[Corpo Rígido com Resistência Linear (Lei de Stokes):] Parâmetros anteriores mais massa ($m$), viscosidade do ar ($\eta$) e raio da esfera ($r$). O coeficiente $c = 6\pi\eta r$ é fixo para cada experimento e representa a intensidade da força de arrasto proporcional à velocidade.\\
    \begin{table}[H]
\centering
\caption{Tabela – Dados referentes a Parametros linear.}
\label{tab:parametros_linear}
\begin{tabular}{llll}
\toprule
Parâmetro & Descrição & Valor & Unidade \\
\midrule
H & Altura Inicial & 10000.0 & m \\
V & Velocidade Inicial & 0.0 & m/s \\
g & Aceleração da Gravidade & 9.81 & m/s² \\
m & Massa & 1.0 & kg \\
eta & Viscosidade do Ar & 0.0 & Pa·s \\
r & Raio da Esfera & 0.01 & m \\
k & Coeficiente Linear (Stokes) & 0.0 & 1/s \\
\bottomrule
\end{tabular}

\end{table}

    \label{tab:parametros_linear}
    \item[Corpo Rígido com Resistência Quadrática:] Parâmetros anteriores mais coeficiente de arrasto ($C_d$), área frontal ($A$) e densidade do ar ($\rho$). O coeficiente $k = \frac{1}{2} \rho C_d A$ é fixo para cada experimento e representa a intensidade da força de arrasto proporcional ao quadrado da velocidade.\\
    \begin{table}[H]
\centering
\caption{Tabela – Dados referentes a Parametros quadratica.}
\label{tab:parametros_quadratica}
\begin{tabular}{llll}
\toprule
Parâmetro & Descrição & Valor & Unidade \\
\midrule
H & Altura Inicial & 10000.0 & m \\
V & Velocidade Inicial & 0.0 & m/s \\
g & Aceleração da Gravidade & 9.81 & m/s² \\
m & Massa & 1.0 & kg \\
Cd & Coeficiente de Arrasto & 0.47 & - \\
A & Área Frontal & 0.01 & m² \\
rho & Densidade do Ar & 1.225 & kg/m³ \\
c & Coeficiente Quadrático & 0.0029 & kg/m \\
\bottomrule
\end{tabular}

\end{table}

    \label{tab:parametros_quadratica}
\end{description}

\textit{Nota:} A notação $c$ (linear) e $k$ (quadrático) é utilizada para diferenciar claramente os regimes de resistência ao longo do texto, experimentos e arquivos de configuração, facilitando a comparação dos efeitos de cada modelo. Os valores adotados para cada cenário estão detalhados nas tabelas acima, o que garante rastreabilidade, comparação entre modelos e reprodutibilidade científica.

\textbf{Parâmetros e análise de acurácia:}
\begin{table}[H]
\centering
\caption{Tabela – Dados referentes a Parâmetros acurácia log.}
\label{tab:parametros_acuracia_log}
\begin{tabular}{l l p{6cm} l}
\toprule
Parâmetro & Descrição & Valor & Unidade \\
\midrule
H & Altura Inicial & 10000.0 & m \\
V & Velocidade Inicial & 0.0 & m/s \\
g & Aceleração da Gravidade & 9.81 & m/s² \\
k & Coeficiente Linear & 3.392920065876977e-06 & 1/s \\
c & Coeficiente Quadrático & 0.00287875 & kg/m \\
dt\_list & Passos de Tempo Testados &
\parbox{6cm}{\scriptsize
1.0000e-04, 3.1623e-04, 1.0000e-03 \\
3.1623e-03, 1.0000e-02, 3.1623e-02 \\
1.0000e-01, 3.1623e-01, 1.0000e+00
}
& s \\
\bottomrule
\end{tabular}
\end{table}

\noindent
Os parâmetros utilizados para o estudo de acurácia, especialmente a lista de passos de tempo testados ($\Delta t$), estão sintetizados na Tabela~\ref{tab:parametros_acuracia_log}. Essa tabela permite avaliar quantitativamente a influência da discretização temporal na precisão dos métodos numéricos.
\label{tab:parametros_acuracia_log}

O erro relativo foi utilizado como métrica principal para avaliar a precisão dos métodos numéricos, sendo calculado por:
\begin{equation}
    E = \left| \frac{v_{\text{num}} - v_{\text{ana}}}{v_{\text{ana}}} \right|
\end{equation}
onde $E$ representa o erro relativo, $v_{\text{num}}$ é o valor numérico obtido na simulação e $v_{\text{ana}}$ o valor analítico de referência.

\textbf{Equações analíticas utilizadas para validação:}
\begin{itemize}
    \item \textbf{Sem resistência do ar:}
    \begin{equation}
        v_{\text{final}} = V - g \cdot t
    \end{equation}
    \item \textbf{Com resistência do ar linear (Lei de Stokes):}
    \begin{equation}
        v_{\text{final}} = \left(V + \frac{mg}{c}\right) e^{-\frac{c}{m} t} - \frac{mg}{c}
    \end{equation}
    \item \textbf{Com resistência do ar quadrática:}
    \begin{equation}
        v_{\text{final}} = -\sqrt{\frac{mg}{k}} \cdot \frac{1 + A e^{2 \sqrt{\frac{gk}{m}} t}}{1 - A e^{2 \sqrt{\frac{gk}{m}} t}}
    \end{equation}
    onde
    \begin{equation}
        A = \frac{V + \sqrt{\frac{mg}{k}}}{V - \sqrt{\frac{mg}{k}}}
    \end{equation}
\end{itemize}
Essas expressões foram implementadas computacionalmente para automatizar o cálculo do erro relativo entre as soluções numérica e analítica, garantindo uma avaliação quantitativa rigorosa da acurácia dos métodos utilizados.

A metodologia desenvolvida permitiu analisar de forma sistemática e reprodutível a queda livre sob diferentes regimes físicos, utilizando métodos numéricos e análises estatísticas para comparar soluções numéricas e analíticas. Os resultados evidenciaram o impacto do passo de tempo e da resistência do ar na precisão das simulações, validando os modelos computacionais e destacando a importância da escolha adequada dos parâmetros para garantir rigor e confiabilidade científica.